\chapter{Suplementos}
\label{cap:suplementos}

\opening{Un suplemento manda sobre el manual básico solo en lo que dice; en todo lo demás, rige el manual básico.}

\section{Suplementos de equipo opcional}
\fuente{POH, sección 9}

La información de los siguientes suplementos aplica solo cuando el equipo correspondiente está instalado. Cada suplemento complementa o reemplaza al manual básico únicamente en las áreas que indica; para todo lo demás, consulte el POH básico.

\begin{itemize}
\item Suplemento 3: versión policial.
\item Suplemento 4: Mariner / Mariner II.
\item Suplemento de aviónica opcional.
\end{itemize}

Existen además suplementos con información exigida por ciertos países: suplemento brasileño, suplemento FATA (Rusia), suplemento IAC AR y suplemento ucraniano.

\section{Versión policial (suplemento 3)}

Debe adjuntarse al POH cuando el helicóptero está configurado con el equipo opcional de versión policial.

\subsection{Limitaciones}

Se permite el VFR nocturno con las luces de aterrizaje \emph{o} el reflector operativos. Las luces de instrumentos y anticolisión deben estar instaladas y operativas.

\subsection{Procedimientos de emergencia, performance y peso y balance}

Sin cambios.

\subsection{Procedimientos normales}

Agregar al punto 10 del prevuelo (costado izquierdo del fuselaje):

\begin{checklist}
\chk{Parlante PA}{Firme}
\chk{Reflector}{Firme}
\end{checklist}

\subsection{Descripción de sistemas}

Puede instalarse el siguiente equipo opcional:

\begin{description}
\item[Célula.] Puertas de cabina con ventanas tipo burbuja, que mejoran la visibilidad hacia abajo y hacia atrás, en lugar de las estándar.
\item[Sistema eléctrico.] Un alternador de 14~V y 70~A reemplaza la unidad estándar de 60~A.
\item[Reflector.] Dos lámparas suspendidas bajo el mentón de la cabina en un poste que atraviesa el lado izquierdo del mentón y el piso. Se controla con una empuñadura e interruptores en el poste, operable por el pasajero o el piloto. Tiene su propio disyuntor de 20~A y relé de potencia, y solo funciona con las luces de aterrizaje apagadas. Un control de fricción está en el poste, justo sobre el piso. Hay dos instalaciones: Collins Dynamics (lámparas halógenas de cuarzo, con un único interruptor de encendido) y Peak Beam Systems (lámparas de arco de xenón, con cuatro interruptores momentáneos que controlan encendido, modo estroboscópico y apertura del haz de 1° (punto) a 40° (difuso)).
\item[Sistema PA.] Un parlante y sirena Whelen de 100~W en el montante trasero izquierdo del tren. El panel de control, en el panel de instrumentos, permite al piloto u observador elegir PA, radio, \emph{yelp} o sirena.
\item[Interruptor remoto de transmisión.] Un interruptor de mano que se enchufa en el receptáculo del frente del asiento izquierdo y permite transmitir al ocupante izquierdo cuando el cíclico izquierdo está retirado.
\item[Panel de disyuntores.] Incluye disyuntores para PA/sirena, ICS, COM~1, transpondedor, COM~2, calefacción, estrobo y luces de navegación.
\end{description}

\begin{precaucion}
Al operar el reflector con el cíclico izquierdo instalado, tenga cuidado de no interferir con el cíclico.
\end{precaucion}

\subsection{Manejo y mantenimiento: retiro del reflector}

\begin{pasos}
\item Afloje, sin retirarlos, los dos pernos de 1/4~in de la abrazadera en el extremo inferior del poste.
\item Retire el pasador de horquilla que une el extremo inferior del tubo central del poste con los brazos entre las lámparas.
\item Desconecte los cables de las lámparas en los conectores de sus carcasas.
\item Tire del anillo del pasador con resorte de la abrazadera del poste y deslice la abrazadera con las lámparas fuera del poste.
\item Retire los cuatro tornillos que fijan la brida del poste al piso. Levante el poste unas pulgadas y desconecte los cables en su conector. Antes de soltarlo, el anillo O del conector debe sacarse de los ganchos de retención. Levante el poste y sáquelo del piso.
\item Cubra con cinta el conector bajo el piso e instale la tapa de chapa suministrada sobre el agujero.
\end{pasos}

La instalación es inversa, excepto: (1) ubique el anillo O del conector bajo los ganchos de retención; (2) la palabra ``FWD'' de la brida debe mirar hacia adelante; (3) asegúrese de que el pasador con resorte encaje en su lugar y apriete los pernos de 1/4~in a 90~in-lb (incluido el torque de autobloqueo de la tuerca).

\section{Mariner y Mariner II (suplemento 4)}

\subsection{Generalidades}

Los R22 Mariner y Mariner II se equipan con tren de flotadores o tren estándar, protección adicional contra la corrosión y un ángulo de incidencia modificado del estabilizador horizontal. Salvo lo indicado en el suplemento, la información del manual básico para el Beta aplica al Mariner, y la del Beta II al Mariner II.

El tren de flotadores consiste en flotadores inflados unidos a los tubos de los patines y extensiones de patín que sostienen la parte trasera de los flotadores. Con flotadores se instala un estabilizador adicional en la base del estabilizador vertical inferior para mejorar la controlabilidad. Se permite cambiar entre tren de flotadores y tren estándar.

\subsection{Limitaciones}

\textbf{\textsc{v}\textsubscript{NE} con flotadores:} con potencia, 95~KIAS hasta 3000~ft de altitud de densidad; sin potencia, 80~KIAS hasta 7500~ft. Sobre esas altitudes, vea los letreros.

\textbf{CG:} vea la figura~\ref{fig:cg-mariner}. Dato 100~in delante del eje del rotor principal.

\textbf{Tipo de operación:} con flotadores instalados, solo VFR diurno.

\widefigure[0.74\linewidth]{assets/poh/mariner-155.png}{Límites de CG del R22 Mariner/Mariner II, con flotadores (línea discontinua) y sin flotadores (línea continua).\label{fig:cg-mariner}}

\widefigure[0.82\linewidth]{assets/poh/mariner-156.png}{Letreros del Mariner/Mariner II: tipo de operación aprobado (día y noche VFR sin flotadores, solo día VFR con flotadores) y tablas de \textsc{v}\textsubscript{NE} con y sin flotadores. Para autorrotación, use 80~KIAS sobre la línea gruesa. Sobre los flotadores: ``ver el manual del piloto para la presión de inflado correcta''.}

\subsection{Límites de presión de los flotadores}

\begin{checklist}
\chk{Presión mínima}{2 psig}
\chk{Presión máxima}{6 psig}
\end{checklist}

El gráfico de la figura~\ref{fig:flotadores} da la presión mínima inicial requerida para una disminución de altitud o de temperatura entre el punto de partida y el destino. Los aumentos de altitud o temperatura pueden ignorarse, porque las válvulas de alivio limitan la presión a 6~psig.

\textbf{Ejemplo:} destino a 1000~ft y 15~°C; condiciones iniciales 5500~ft y 5~°C. El cambio esperado es $-4500$~ft y $+10$~°C. En el gráfico, ubique la línea de $-4500$~ft, lea horizontalmente hasta la línea de $+10$~°C y baje: la presión inicial mínima requerida es de aproximadamente 3,7~psig.

\widefigure[0.74\linewidth]{assets/poh/mariner-157.png}{Presión requerida de los flotadores en función del cambio de altitud y temperatura.\label{fig:flotadores}}

\begin{precaucion}
No mantener la presión requerida puede producir una flotabilidad inadecuada o inestabilidad en vuelo.
\end{precaucion}

\subsection{Procedimientos de emergencia}

Sin flotadores aplican los procedimientos del manual básico. Con flotadores:

\textbf{Falla de potencia sobre 500~ft AGL.} Autorrotación a tierra: igual que el R22 sin flotadores. Autorrotación al agua: pasos 1 a 7 iguales; luego (8) a unos 8~ft AGL aplique cíclico adelante para nivelar y suba el colectivo justo antes del contacto, tocando con actitud levemente de nariz arriba y la nariz recta; (9) mantenga el cíclico en la posición de contacto y el colectivo arriba hasta que se detenga el movimiento de avance.

\textbf{Falla de potencia entre 8 y 500~ft AGL.} Autorrotación a tierra: igual. Autorrotación al agua: pasos 1 a 4 iguales; luego (5) a unos 8~ft AGL nivele y suba el colectivo justo antes del contacto, con actitud levemente de nariz arriba y nariz recta; (6) mantenga el cíclico en la posición de contacto y el colectivo arriba hasta que se detenga el avance.

\begin{precaucion}
Bajar el colectivo o aplicar cíclico adelante mientras el helicóptero avanza sobre el agua puede hundir los flotadores y hacer que el helicóptero capote.
\end{precaucion}

\textbf{Amaraje:} no aplica con flotadores instalados.

\subsection{Procedimientos normales}

Agregar al prevuelo:

\begin{checklist}
\chkhead{14. Flotadores inflables}
\chk{Presión de los flotadores}{Revisar (ver Limitaciones)}
\chk{Estado de los flotadores}{Revisar}
\end{checklist}

\begin{precaucion}
Los helicópteros con flotadores inflados tienen una característica de alabeo adversa: al derrapar con la nariz a la izquierda o a la derecha, el helicóptero tiende a alabear en sentido contrario. Para evitarlo, mantenga el helicóptero compensado con derrape cero y extreme las precauciones al simular fallas de potencia.
\end{precaucion}

\textbf{Operación sobre el agua.} La velocidad máxima recomendada de rodaje sobre el agua es de 5~nudos; se requiere algo de colectivo.

\begin{precaucion}
Si arranca o detiene el rotor sobre el agua, asegúrese de que el área esté despejada: el helicóptero puede girar una vuelta completa mientras la RPM del rotor de cola es baja.
\end{precaucion}

\begin{precaucion}
Para no tocar los topes de balanceo del rotor, no aplique cíclico a fondo durante el rodaje sobre el agua.
\end{precaucion}

\begin{nota}
Se ha demostrado una operación segura sobre el agua con olas de hasta 12~in (de valle a cresta).
\end{nota}

\textbf{Autorrotación de práctica con contacto con el suelo:} igual que sin flotadores, salvo que la parte trasera de los flotadores puede rozar el pavimento. Aplique cinta protectora bajo los flotadores e inspecciónela por desgaste después de los contactos.

\textbf{Autorrotación de práctica con contacto con el agua:} igual que con contacto con el suelo, pero tocando con actitud levemente de nariz arriba y nariz recta, manteniendo el cíclico en la posición de contacto y el colectivo arriba hasta que se detenga el avance (vea la precaución anterior sobre capotaje).

\subsection{Performance}

\widefigure[0.74\linewidth]{assets/poh/mariner-161.png}{Techo de vuelo estacionario IGE vs. peso bruto --- R22 Mariner (O-320-B2C), 2~ft de separación, acelerador a fondo y 104\,\% RPM.}

\widefigure[0.74\linewidth]{assets/poh/mariner-162.png}{Techo de vuelo estacionario OGE vs. peso bruto --- R22 Mariner (O-320-B2C), potencia de despegue de 5 minutos o acelerador a fondo, 104\,\% RPM.}

La performance en estacionario del Mariner II equivale a la del Beta II.

\subsection{Peso y balance}

\begin{precaucion}
Al cambiar de configuración con flotadores a sin flotadores, o viceversa, debe revisarse el peso y balance del helicóptero básico y reajustarse la RPM de autorrotación según el Manual de Mantenimiento.
\end{precaucion}

El peso vacío básico y el CG en ambas configuraciones se incluyen en los datos de peso y balance entregados con el helicóptero; las modificaciones se registran en el registro del manual básico. Determine los pesos y momentos del helicóptero cargado según las instrucciones de carga; deben quedar dentro de la envolvente de la figura~\ref{fig:momento-mariner}.

\widefigure[0.74\linewidth]{assets/poh/mariner-163.png}{Envolvente de momento del CG admisible vs. peso bruto, R22 Mariner/Mariner II (sin flotadores: línea continua; con flotadores: discontinua).\label{fig:momento-mariner}}

\subsection{Manejo, servicio y mantenimiento}

Con flotadores se requieren ruedas de maniobra especiales (vea el Manual de Mantenimiento). Para prolongar la vida de los flotadores:

\begin{pasos}
\item No los infle a una presión mayor que la requerida en Limitaciones.
\item No los infle arbitrariamente hasta la presión de la válvula de alivio.
\item Reduzca la presión si el calentamiento solar causa una presión excesiva.
\item No deje los flotadores desinflados por largos períodos; mantenga siempre algo de presión para que conserven su forma.
\end{pasos}

\begin{precaucion}
Al inflar las cámaras individualmente (sin múltiple), aumente la presión de cada cámara en incrementos de 0,5~psig.
\end{precaucion}

\subsection{Consejos de seguridad}

Las características de vuelo y manejo de un helicóptero con flotadores inflados son más críticas que con tren convencional. Por ejemplo, tiene una característica de alabeo adversa: al derrapar con la nariz a la derecha o a la izquierda, tiende a alabear en sentido contrario, hacia afuera del viraje. Esto podría ser extremadamente peligroso si un alumno o instructor no aplica pedal derecho, o aplica el pedal equivocado, durante una falla de potencia simulada. Además, la sustentación aerodinámica de los flotadores dificulta el control de la RPM y de la actitud longitudinal en las entradas en autorrotación, y el helicóptero es más sensible a las ráfagas y más difícil de volar en turbulencia.

Por estas razones se recomienda encarecidamente retirar los flotadores e instalar el tren estándar cuando el helicóptero se use para instrucción de vuelo primaria. Con flotadores instalados, los pilotos deben mantener el helicóptero compensado con derrape cero y extremar las precauciones al simular fallas de potencia.

\section{Aviónica opcional}

Este suplemento debe incluirse en el POH cuando hay instalada cierta aviónica opcional de fábrica. Con cada helicóptero nuevo se entrega un juego de instrucciones del fabricante de toda la aviónica instalada. El suplemento cubre el \textbf{Aspen Avionics EFD 1000H PFD}.

\begin{nota}
En la instalación de pantalla primaria de vuelo (PFD) de Robinson se conservan el anemómetro, el altímetro, la brújula, el tacómetro y los instrumentos de motor. Los pilotos deben usar los instrumentos tradicionales como primarios salvo que estén completamente familiarizados con la aviónica instalada.
\end{nota}

Secciones 2 a 6 y 8: sin cambios.

\textbf{Descripción de sistemas.} El Aspen EFD 1000H es una PFD optimizada para helicópteros, disponible en versión ``Pilot'' (básica) o ``Pro'' (con funciones de navegación más avanzadas). Su documento de referencia es la \emph{Aspen Avionics Evolution Flight Display EFD 1000H PFD Pilot's Guide}, N.º 091-00012-001.

\begin{nota}
Puede usarse un filtro de luz Robinson P/N D327-4 para reducir los reflejos en el parabrisas de noche. Se instala sujetándolo al frente de la pantalla; su uso queda a criterio del piloto.
\end{nota}

\widefigure[0.80\linewidth]{assets/poh/sis-171.png}{Panel de instrumentos opcional con PFD Aspen EFD 1000H (la configuración exacta varía con el equipo opcional y la fecha de fabricación).}
